\section*{अध्याय \ref{ch:b3:point-groups} --- सममिति और बिंदु समूह}

\begin{solution}{exo:b3:point-groups:1}
एथीन: तीन परस्पर लंबवत $C_2$, तीन तल, $i$: $D_{2h}$। \ce{XeF4} (वर्ग
समतलीय): $C_4$, चार $C_2 \perp C_4$, $\sigma_h$, $2\sigma_v$, $2\sigma_d$, $i$, $S_4$:
$D_{4h}$। \ce{PCl5} (त्रिकोणीय द्विपिरैमिड): $C_3$, तीन $C_2$, $\sigma_h$, $3\sigma_v$,
$S_3$: $D_{3h}$। 1,3,5-ट्राइक्लोरोबेंज़ीन: \ce{BF3} वाले ही तत्व: $D_{3h}$।
\end{solution}

\begin{solution}{exo:b3:point-groups:2}
सांतरित एथेन: $C_3$, तीन लंबवत $C_2$, तीन $\sigma_d$, $i$, $S_6$:
$D_{3d}$। ग्रसित: $C_3$, तीन $C_2$, $\sigma_h$, तीन $\sigma_v$: $D_{3h}$।
कुर्सी साइक्लोहेक्सेन: $D_{3d}$।
\end{solution}

\begin{solution}{exo:b3:point-groups:3}
\ce{SO3} ($D_{3h}$): नहीं। \ce{SO2} ($C_{2v}$): हाँ। \ce{PCl3} ($C_{3v}$): हाँ।
\ce{XeF4} ($D_{4h}$): नहीं। \ce{CH2Cl2} ($C_{2v}$): हाँ। cis-1,2-डाइक्लोरोएथीन
($C_{2v}$): हाँ; trans ($C_{2h}$): नहीं।
\end{solution}

\begin{solution}{exo:b3:point-groups:4}
$C_2(z) = \mathrm{diag}(-1,-1,1)$, $\chi = -1$; $i = \mathrm{diag}(-1,-1,-1)$,
$\chi = -3$; $\sigma_h(xy) = \mathrm{diag}(1,1,-1)$, $\chi = 1$।
\end{solution}

\begin{solution}{exo:b3:point-groups:5}
पिछले अभ्यास के विकर्ण आव्यूहों से: $C_2i = \sigma_h$, $C_2\sigma_h =
i$, $i\sigma_h = C_2$, हर संक्रिया स्वयं अपना प्रतिलोम है, और हर गुणनफल
क्रमविनिमेय है। समूह क्रमविनिमेय है, अतः $X^{-1}AX = A$, सभी $X$ के लिए: हर
संक्रिया अकेली एक वर्ग है (चार वर्ग, विमा 1 के चार \omterm{def:b3:point-groups:irrep}{अलघुकरणीय
निरूपण})।
\end{solution}

\begin{solution}{exo:b3:point-groups:6}
$\sigma_v(1)$ को $xz$ तल लीजिए। बिंदु $(1,0)$, $C_3$ द्वारा $(-\frac12,
\frac{\sqrt3}2)$ पर जाता है, फिर $\sigma_v(1)$ द्वारा $(-\frac12,-\frac{\sqrt3}2)$ पर; दूसरे
क्रम में वह $(1,0)$ पर, फिर $(-\frac12,\frac{\sqrt3}2)$ पर जाता है। दोनों गुणनफल
दो भिन्न तलों ($\sigma_v(3)$ और $\sigma_v(2)$) में परावर्तन हैं: समूह
क्रमविनिमेय नहीं है, इसीलिए उसके परावर्तन तीन का एक वर्ग बनाते हैं।
\end{solution}

\begin{solution}{exo:b3:point-groups:7}
ऐंठा हुआ बाइफ़ेनिल वलयों के बीच के आबंध के अनुदिश $C_2$ और उसके लंबवत दो $C_2$
बनाए रखता है, जो वलय-तलों के बीच के कोणों को समद्विभाजित करते हैं; समतल या लंबवत
रूपों के सभी तल और \omterm{def:b3:point-groups:inversion}{प्रतिलोमन केंद्र}
खो जाते हैं। तीन परस्पर लंबवत $C_2$ और किसी अनुचित तत्व के बिना समूह
$D_2$ है: अणु किरैल है (उसके दोनों ऐंठे रूप प्रतिबिंबरूपी हैं, जो
पृथक किए जा सकते हैं जब भारी ऑर्थो प्रतिस्थापी घूर्णन को रोकते हैं)।
\end{solution}

\begin{solution}{exo:b3:point-groups:8}
$d_{z^2}$ और $d_{x^2-y^2}$: $a_1$ ($x^2$, $y^2$, $z^2$, $A_1$ हैं); $d_{xy}$: $a_2$; $d_{xz}$:
$b_1$; $d_{yz}$: $b_2$।
\end{solution}

\begin{solution}{exo:b3:point-groups:9}
$1 + 1 + 4 + 9 + 9 = 24 = h$। $\sum g\,\chi_{A_1}\chi_{T_2} = 1\cdot3 + 8\cdot0 + 3(-1)
+ 6(-1) + 6(1) = 0$।
\end{solution}

\begin{solution}{exo:b3:point-groups:10}
प्रतिच्छेद रेखा के लंबवत तल में, कोण
$\alpha$ पर स्थित रेखा में परावर्तन ध्रुवी कोण $\phi$ को $2\alpha - \phi$ पर ले जाता है। दो परावर्तन,
पहले $\alpha$ पर फिर $\alpha + \theta$ पर: $\phi \to 2\alpha - \phi \to 2(\alpha + \theta) - (2\alpha - \phi) =
\phi + 2\theta$, रेखा के परितः $2\theta$ का घूर्णन। इस प्रकार
$60^\circ$ पर स्थित दो ऊर्ध्वाधर तल $120^\circ$ का घूर्णन उत्पन्न करते हैं: एक $C_3$।
\end{solution}

\begin{solution}{exo:b3:point-groups:11}
C=C=C अक्ष एक $C_2$ और एक $S_4$ है ($90^\circ$ घुमाइए, जिससे
दोनों \ce{CH2} तल परस्पर बदल जाते हैं, फिर केंद्रीय कार्बन के तल में परावर्तन कीजिए); दोनों
\ce{CH2} तल $\sigma_d$ हैं; अक्ष के लंबवत दो $C_2$ उन्हें समद्विभाजित करते हैं।
कोटि 8: $D_{2d}$, अकिरैल। \ce{ClHC=C=CHCl} में प्रतिस्थापन से तल, $S_4$ और अक्षीय
$C_2$ नष्ट हो जाते हैं; अक्ष के लंबवत केवल एक $C_2$
बचता है: $C_2$, किरैल --- एक अक्षीय रूप से किरैल ऐलीन।
\end{solution}

\begin{solution}{exo:b3:point-groups:12}
\ce{SF6}: $O_h$ ($h = 48$)। \ce{SF5Cl}: Cl से होकर $C_4$ और चार $\sigma_v$:
$C_{4v}$ ($h = 8$)। trans-\ce{SF4Cl2}: $D_{4h}$ ($h = 16$)। cis-\ce{SF4Cl2}: $C_{2v}$ ($h
= 4$)। हर एक $O_h$ की संक्रियाओं का एक उपसमुच्चय रखता है जो गुणनफलों के अंतर्गत संवृत है,
और 8, 16, 4, 48 को विभाजित करते हैं। ध्रुवीय: \ce{SF5Cl} और cis-\ce{SF4Cl2}।
\end{solution}

\begin{solution}{pb:b3:point-groups:1}
\textbf{1.} C और हर H से होकर जाने वाले चार काय-विकर्णों के अनुदिश; हर एक
$C_3$ और $C_3^2$ देता है: 8 घूर्णन।
\textbf{2.} फलक-केंद्रों से होकर $x$, $y$, $z$ के अनुदिश: $C_2(z)$, $(1,1,1)
\to (-1,-1,1)$ करता है, जो एक हाइड्रोजन है। $90^\circ$ का घूर्णन, $z$ के परितः, और उसके बाद
$xy$ में परावर्तन, $(1,1,1) \to (-1,1,1) \to (-1,1,-1)$ करता है, एक हाइड्रोजन: एक
$S_4$; $S_4^3$ सहित, 6 संक्रियाएँ।
\textbf{3.} तल $x = \pm y$, $y = \pm z$, $x = \pm z$; $x = y$ में
$(1,1,1)$ और $(-1,-1,1)$ समाहित हैं: हर तल में दो C--H आबंध हैं।
\textbf{4.} $1 + 8 + 3 + 6 + 6 = 24$।
\textbf{5.} प्रतिलोमन $(1,1,1)$ को $(-1,-1,-1)$ पर ले जाता है, जो हाइड्रोजन नहीं है:
कोई $i$ नहीं।
\textbf{6.} $E$; $8C_3$; $3C_2$; $6S_4$; $6\sigma_d$।
\textbf{7.} \ce{CH3Cl}: $C_{3v}$, $h = 6$।
\textbf{8.} \ce{CH2Cl2}: $C_{2v}$, $h = 4$।
\textbf{9.} \ce{CHCl3}: $C_{3v}$; \ce{CCl4}: $T_d$, $h = 24$।
\textbf{10.} \ce{CH2FCl}: $C_s$ ($h = 2$); \ce{CHFClBr}: $C_1$ ($h = 1$)।
\textbf{11.} हर एक $T_d$ की वे संक्रियाएँ रखता है जो
प्रतिस्थापन का सम्मान करती हैं; 6, 4, 6, 24, 2, 1 सभी 24 को विभाजित करते हैं।
\textbf{12.} तीनों हाइड्रोजन, जिन्हें $C_3$, $C_3^2$ और तीनों
$\sigma_v$ क्रमचयित करते हैं।
\textbf{13.} \ce{CCl4} को छोड़कर सभी: \ce{CH3Cl} और \ce{CHCl3} के लिए $C_3$ अक्ष के अनुदिश,
\ce{CH2Cl2} के लिए $C_2$ अक्ष के अनुदिश, \ce{CH2FCl} के लिए \omterm{def:b3:point-groups:mirror}{दर्पण तल} में, \ce{CHFClBr} के लिए
किसी आरोपित दिशा में नहीं।
\textbf{14.} केवल \ce{CHFClBr} ($C_1$, कोई $S_n$ नहीं)।
\textbf{15.} दो प्रतिबिंबरूपी।
\textbf{16.} उसके दोनों हाइड्रोजन तुल्य हैं: C, F और Cl से होकर जाने वाला, H--C--H को
समद्विभाजित करने वाला तल एक \omterm{def:b3:point-groups:mirror}{दर्पण तल} है।
\textbf{17.} नहीं: चार भिन्न प्रतिस्थापी केवल $E$ छोड़ते हैं; किसी $S_n$ के बिना
अणु किरैल है (और ध्रुवीय हो सकता है)।
\textbf{18.} $(x, y, z)$: $T_2$।
\textbf{19.} दो हाइड्रोजन चतुष्फलक के एक किनारे के सिरे हैं; 24
संक्रियाएँ किसी भी किनारे को किसी अन्य पर ले जाती हैं (6 किनारे एक ही समुच्चय बनाते हैं): एक
ही \ce{CH2Cl2}।
\textbf{20.} \ce{CH2FCl}: एक। \ce{CHFClBr}: दो, प्रतिबिंबरूपियों का एक युग्म।
\textbf{21.} तीन: ऑर्थो ($C_{2v}$), मेटा ($C_{2v}$), पैरा ($D_{2h}$)।
\textbf{22.} तीन: 1,2,3- ($C_{2v}$), 1,2,4- ($C_s$), 1,3,5- ($D_{3h}$)।
\textbf{23.} दो (Cl परमाणु सन्निकट या सम्मुख)। केवल एक \ce{CH2Cl2}
ज्ञात है, जिसने समतल कार्बन को अस्वीकार किया और 1874 में प्रस्तावित
चतुष्फलकीय कार्बन का समर्थन किया।
\textbf{24.} $1^2 + 1^2 + 2^2 + 3^2 + 3^2 = 24$: \textbf{मेथेन के \omterm{def:b3:point-groups:point-group}{बिंदु समूह} $T_d$ की
कोटि $h = 24$ है}।
\end{solution}
